%! The Tangent Function — Unit 7, Lesson 7.1 — Lesson Plan
\documentclass[10pt]{article}
\usepackage{precalculus-article}
\usepackage{precalculus-boxes}
\usepackage{graphicx}
\graphicspath{{images/}}
\newcommand{\TallMath}[1]{$\displaystyle #1\rule[-1.4em]{0pt}{3.2em}$}
\newcommand{\CourseName}{Precalculus}
\newcommand{\MeetingLength}{60 minutes}
\newcommand{\UnitNumberName}{Unit 7: Analytic Trigonometry and Polar Coordinates \quad}
\newcommand{\LessonNumberName}{Lesson 7.1: The Tangent Function}

\begin{document}

%-----------------------------------------------------------------------
% 1. TITLE BLOCK
%-----------------------------------------------------------------------
\begin{center}
  {\Large\bfseries\color{navy} Precalculus}\\[4pt]
  {\large\bfseries Unit 7: Analytic Trigonometry and Polar Coordinates}\\[2pt]
  {\large\bfseries Lesson 7.1: The Tangent Function}\\[2pt]
  {\normalsize (2 instructional periods, \MeetingLength\ each)}
\end{center}

\vspace{4pt}
\hrule
\vspace{8pt}

%-----------------------------------------------------------------------
% 2. PRIMARY OBJECTIVE
%-----------------------------------------------------------------------
\begin{tcolorbox}[colback=sky, colframe=navy, boxrule=0.9pt, arc=2mm,
    left=2mm, right=2mm, top=1.5mm, bottom=1.5mm]
  \textbf{Primary Objective} \hfill \textit{Big Idea: Functions}

  \medskip
  Students construct representations of the tangent function using the unit
  circle, describe its key characteristics (period $\pi$, vertical
  asymptotes, monotone increase, concavity change), and analyze the effect
  of additive and multiplicative transformations on the tangent graph.
\end{tcolorbox}

\vspace{6pt}

%-----------------------------------------------------------------------
% 3. PRIORITY IDEAS & SKILLS
%-----------------------------------------------------------------------
\begin{skillbox}[Priority Ideas \& Skills]{goldbox}
  \textbf{Skills:} 2.A (Identify information from representations) \quad
  3.A (Describe characteristics of functions)

  \medskip
  \textbf{Essential Knowledge:}
  \begin{itemize}
    \item \textbf{EK 3.8.A.1/A.2} $f(\theta)=\tan\theta$ gives the slope of the
          terminal ray; $\tan\theta = \sin\theta/\cos\theta$ where $\cos\theta\neq 0$.
    \item \textbf{EK 3.8.B.1} Period is $\pi$ --- slope of the terminal ray repeats
          every half-revolution.
    \item \textbf{EK 3.8.B.2} Vertical asymptotes at $\theta = \pi/2 + k\pi$,
          $k\in\mathbb{Z}$, because $\cos\theta=0$ there.
    \item \textbf{EK 3.8.B.3} The tangent function is always increasing between
          consecutive asymptotes; its graph changes from concave down to concave
          up within each period.
    \item \textbf{EK 3.8.C.1--C.5} Transformations $g(\theta) = a\cdot\tan(b(\theta+c))+d$:
          vertical dilation $|a|$, period $\pi/|b|$, phase shift $-c$, vertical
          translation $d$.
  \end{itemize}
\end{skillbox}

%-----------------------------------------------------------------------
% 4. VOCABULARY
%-----------------------------------------------------------------------
\begin{skillbox}[{Vocabulary, Concepts, \& Theorems}]{greenbox}
  \begin{multicols}{3}
  \begin{itemize}
    \item tangent function
    \item slope of terminal ray
    \item vertical asymptote
    \item period (tangent)
    \item phase shift
    \item vertical dilation
    \item concavity
    \item inflection point
  \end{itemize}
  \end{multicols}
\end{skillbox}

%-----------------------------------------------------------------------
% 5. ACTIVATE PRIOR KNOWLEDGE
%-----------------------------------------------------------------------
\begin{skillbox}[Activate Prior Knowledge \& Spiral Review]{sky}
  \begin{itemize}
    \item \textbf{Lessons 3.2--3.4} --- Unit circle: $\sin\theta$ and
          $\cos\theta$ values at multiples of $\pi/6$ and $\pi/4$; interpreting
          a point on the unit circle as $(\cos\theta, \sin\theta)$.
    \item \textbf{Algebra} --- Slope of a line through the origin: rise over run.
    \item \textbf{Pre-3.8} --- Identifying vertical asymptotes of rational
          functions; behavior near a denominator zero.
  \end{itemize}
\end{skillbox}

%-----------------------------------------------------------------------
% 6. HOOK
%-----------------------------------------------------------------------
\begin{skillbox}[Hook]{sky}
  \textit{``A lighthouse beam sweeps the shore. A sensor at the end of the
  dock records how far along the shore the beam strikes as a function of the
  angle $\theta$ the beam makes with the perpendicular. When the beam is
  nearly parallel to the shore, the illuminated point races off to infinity.
  What function models this?''}

  \medskip
  \begin{itemize}
    \item If the dock is 1 unit from the lighthouse, the distance along the
          shore is $\tan\theta$.
    \item As $\theta\to\pi/2$, $\tan\theta\to+\infty$: the beam sweeps to
          infinity --- an asymptote in action.
    \item \textbf{Key idea:} the tangent function is the slope of the terminal
          ray of angle $\theta$ on the unit circle.
  \end{itemize}

  \smallskip
  \textit{Transition:} Today we derive $\tan\theta = \sin\theta/\cos\theta$
  from the unit circle, describe all key characteristics of the tangent graph,
  and explore how transformations shift and stretch it.
\end{skillbox}

%-----------------------------------------------------------------------
% 7. LESSON
%-----------------------------------------------------------------------
\begin{skillbox}[Lesson --- Period 1: Deriving and Graphing Tangent]{sky}
  \begin{multicols}{2}
  \textbf{Tangent as Slope of Terminal Ray}

  Place angle $\theta$ in standard position. The terminal ray passes
  through $(\cos\theta,\sin\theta)$ and the origin. Its slope is:
  \[
    \tan\theta = \frac{\sin\theta}{\cos\theta}
  \]
  Defined whenever $\cos\theta\neq 0$.

  \columnbreak

  \textbf{Key Values Table}

  \renewcommand{\arraystretch}{1.35}
  \begin{tabular}{|c|c|}
  \hline
  $\theta$ & $\tan\theta$ \\
  \hline
  $0$        & $0$ \\
  $\pi/4$    & $1$ \\
  $\pi/3$    & $\sqrt{3}$ \\
  $\pi/2$    & undefined \\
  $2\pi/3$   & $-\sqrt{3}$ \\
  $3\pi/4$   & $-1$ \\
  $\pi$      & $0$ \\
  \hline
  \end{tabular}
  \end{multicols}

  \textbf{Why period $= \pi$:} After a half-revolution the terminal ray
  points in the opposite direction, but its slope is identical: the rise
  and run both change sign, so the ratio is unchanged.  Thus
  $\tan(\theta+\pi)=\tan\theta$.

  \medskip
  \textbf{Asymptotes:} When $\cos\theta=0$ (i.e.\ $\theta=\pi/2+k\pi$),
  the slope is undefined and $|\tan\theta|$ grows without bound --- vertical
  asymptotes at each of these values.
\end{skillbox}

\begin{skillbox}[Lesson (cont.) --- Period 2: Characteristics \& Transformations]{sky}
  \begin{multicols}{2}
  \textbf{Characteristics in $(-\pi/2,\,\pi/2)$}
  \begin{itemize}
    \item Passes through $(0,0)$; always increasing.
    \item Concave down for $\theta<0$; inflection at $(0,0)$;
          concave up for $\theta>0$.
    \item $\tan\theta\to+\infty$ as $\theta\to(\pi/2)^-$.
    \item $\tan\theta\to-\infty$ as $\theta\to(-\pi/2)^+$.
  \end{itemize}

  \columnbreak

  \textbf{Transformations: $g(\theta)=a\tan(b(\theta+c))+d$}

  \renewcommand{\arraystretch}{1.4}
  \begin{tabular}{@{}ll@{}}
  Parameter & Effect \\
  \hline
  $a$ & vertical dilation by $|a|$ \\
  $b$ & period becomes $\pi/|b|$ \\
  $c$ & phase shift of $-c$ \\
  $d$ & vertical translation by $d$ \\
  \end{tabular}
  \end{multicols}

  \textbf{Finding new asymptotes:} Set $b(\theta+c)=\pi/2+k\pi$ and
  solve for $\theta$. Example: for $g(\theta)=\tan(2\theta)$,
  asymptotes where $2\theta=\pi/2+k\pi$, so $\theta=\pi/4+k\pi/2$.
\end{skillbox}

%-----------------------------------------------------------------------
% 8. EXPLICIT INSTRUCTION
%-----------------------------------------------------------------------
\begin{skillbox}[Explicit Instruction: Describing Tangent Transformations]{sky}
  \textbf{Steps:}
  \begin{enumerate}
    \item Identify $a$, $b$, $c$, $d$ from $g(\theta)=a\tan(b(\theta+c))+d$.
    \item State period: $\pi/|b|$.
    \item Find asymptotes: solve $b(\theta+c)=\pm\pi/2$.
    \item State phase shift $-c$ and vertical shift $d$.
    \item State vertical dilation: output values scaled by $|a|$.
  \end{enumerate}

  \textbf{Worked Example:} $g(\theta)=2\tan\!\left(\tfrac{1}{2}(\theta-\tfrac{\pi}{4})\right)+1$

  \begin{enumerate}
    \item $a=2$, $b=1/2$, $c=-\pi/4$, $d=1$.
    \item Period $= \pi\div(1/2) = 2\pi$.
    \item Asymptotes: $\tfrac{1}{2}(\theta-\tfrac{\pi}{4})=\pm\tfrac{\pi}{2}$
          $\Rightarrow$ $\theta=\tfrac{\pi}{4}\pm\pi$, so
          $\theta=\tfrac{5\pi}{4}$ and $\theta=-\tfrac{3\pi}{4}$
          (plus integer multiples of $2\pi$).
    \item Phase shift: $+\pi/4$ (right); vertical shift: up $1$.
    \item All output values doubled (vertical dilation by 2).
  \end{enumerate}
\end{skillbox}

%-----------------------------------------------------------------------
% 9. ACTIVE MONITORING
%-----------------------------------------------------------------------
\begin{skillbox}[Active Monitoring]{redbox}
  Circulate as students work. Watch for:
  \begin{itemize}
    \item Confusing the period of tangent ($\pi$) with that of sine/cosine ($2\pi$).
    \item Placing asymptotes at multiples of $\pi$ instead of $\pi/2+k\pi$.
    \item Treating $|a|$ as an ``amplitude'' (tangent has no amplitude).
    \item Forgetting to divide by $|b|$ when computing the new period.
  \end{itemize}

  \textbf{Cold-call:} ``What is the period of $\tan(3\theta)$?
  Where are its asymptotes?''
  \quad (Period $= \pi/3$; asymptotes at $\theta = \pi/6 + k\pi/3$.)
\end{skillbox}

%-----------------------------------------------------------------------
% 10. GROUP WORK & DIFFERENTIATION
%-----------------------------------------------------------------------
\begin{skillbox}[Group Work \& Differentiation]{redbox}
  Students work in assigned tiers on the Group Activity.
  \begin{itemize}
    \item \textbf{Tier R (Remediate):} Given unit-circle $\sin/\cos$ values,
          compute $\tan\theta$ at key angles and identify asymptote locations.
    \item \textbf{Tier A (Approaching Proficiency):} Given a tangent graph,
          identify period, asymptotes, and concavity; describe the
          transformation $g(\theta)=a\tan(\theta)+d$.
    \item \textbf{Tier E (Extension):} Analyze
          $g(\theta)=2\tan(\theta/2-\pi/4)+1$: identify all parameters,
          describe the full transformation, find asymptotes, and evaluate $g$
          at two specific angles.
  \end{itemize}
\end{skillbox}

%-----------------------------------------------------------------------
% 11. INDIVIDUAL WORK & ASSESSMENT
%-----------------------------------------------------------------------
\begin{skillbox}[Individual Work \& Assessment]{redbox}
  \begin{itemize}
    \item \textbf{Exit Ticket} (last 5--7 min): 3 short items --- computing
          $\tan\theta$ from given $\sin/\cos$ values, identifying the period
          of a transformed tangent, and locating one asymptote.
    \item \textbf{Homework}: 5--6 problems including an AP-style MC and a
          preview of Lesson 3.9 (inverse trigonometric functions).
    \item \textbf{Formative data:} Collect exit tickets; students who
          misplace asymptotes need reteaching before Lesson 3.9.
  \end{itemize}
\end{skillbox}

%-----------------------------------------------------------------------
% 12. REINFORCEMENT & EXTENSION
%-----------------------------------------------------------------------
\begin{skillbox}[Reinforcement \& Extension]{goldbox}
  \textbf{Reinforce:} Return to the unit circle. Draw several terminal
  rays and compute their slopes directly.  The slope of the ray through the
  origin and $(\cos\theta,\sin\theta)$ is $\sin\theta/\cos\theta$ --- tangent
  is slope.

  \medskip
  \textbf{Extension:} Use Desmos to graph $y=a\tan(\theta)$ for several
  values of $a$. Ask: as $|a|$ increases, how does the graph steepen near
  its inflection point?

  \medskip
  \textbf{Preview --- Lesson 3.9: Inverse Trigonometric Functions}

  In Lesson 3.9 students restrict the domains of sine, cosine, and tangent
  to make them invertible, define $\arcsin$, $\arccos$, and $\arctan$, and
  evaluate inverse trig expressions --- building directly on today's
  asymptote behavior of $\tan\theta$ near $\pm\pi/2$.
\end{skillbox}

\end{document}
